\section*{第 \ref{ch:b1:series} 章——数值级数}

\begin{solution}{exo:b1:series:1}
$\dfrac{1}{n(n+1)} = \dfrac1n - \dfrac{1}{n+1}$：裂项相消，$S_N =
1 - \frac{1}{N+1} \to 1$。收敛，和为 $1$。

$\ln\bigl(1 - \frac{1}{n^2}\bigr) = \ln\frac{(n-1)(n+1)}{n^2} =
\ln\frac{n-1}{n} - \ln\frac{n}{n+1}$：仍是裂项相消，$S_N =
\ln\frac12 - \ln\frac{N}{N+1} \to -\ln 2$。收敛，和为 $-\ln 2$。

$\dfrac{3^n + 4^n}{5^n} = \bigl(\frac35\bigr)^n +
\bigl(\frac45\bigr)^n$：两个收敛的\omterm{ex:b1:series:geometric}{等比级数}，和为
$\frac{1}{1 - 3/5} + \frac{1}{1 - 4/5} = \frac52 + 5 = \frac{15}{2}$。
\end{solution}

\begin{solution}{exo:b1:series:2}
全部使用\omterm{thm:b1:series:ratio}{比值判别法}。

$\frac{u_{n+1}}{u_n} = \frac{(n+1)^2}{2n^2} \to \frac12 < 1$：
收敛。

$\frac{u_{n+1}}{u_n} = \frac{(n+1)!\,n^n}{n!\,(n+1)^{n+1}} =
\bigl(\frac{n}{n+1}\bigr)^n = \bigl(1 + \frac1n\bigr)^{-n} \to
\frac1\eu < 1$：收敛。

乘上因子 $2^n$ 后，比值 $\to \frac2\eu < 1$：收敛。

取 $3^n$ 时，比值 $\to \frac3\eu > 1$：发散（项趋于
$+\infty$）。
\end{solution}

\begin{solution}{exo:b1:series:3}
各项均非负，利用\omterm{def:b1:logic:equiv}{等价关系}
（见\cref{thm:b1:series:positive}）。

$\sin\frac{1}{n^2} \sim \frac{1}{n^2}$：收敛（黎曼 $\alpha
= 2$）。

$1 - \cos\frac1n \sim \frac{1}{2n^2}$：收敛。

$\frac{1}{\sqrt{n(n+1)}} \sim \frac1n$：发散。

$\frac{\ln n}{n^2} = \frac{1}{n^{3/2}}\cdot\frac{\ln n}{n^{1/2}}$
而 $\frac{\ln n}{\sqrt n} \to 0$
（见\cref{prop:b1:functions:powerrules}），于是 $\frac{\ln n}{n^2} \leq
\frac{1}{n^{3/2}}$（$n$ 充分大时）：收敛。
\end{solution}

\begin{solution}{exo:b1:series:4}
当 $n \geq 2$ 时，$\frac{1}{n^2} \leq \frac{1}{n(n-1)} =
\frac{1}{n-1} - \frac1n$。因此
\[
\sum_{n=1}^{N} \frac{1}{n^2} \leq 1 + \sum_{n=2}^{N}
\Bigl(\frac{1}{n-1} - \frac1n\Bigr) = 1 + 1 - \frac1N < 2 :
\]
部分和单调递增且以 $2$ 为\omterm{def:b1:reals:bounds}{上界}，由
\cref{thm:b1:series:positive}可知\omterm{def:b1:series:def}{级数}收敛，和 $\leq 2$。（精确值
$\frac{\pi^2}{6}$ 留待本科二年级再来欣赏。）
\end{solution}

\begin{solution}{exo:b1:series:5}
$\sum \frac{(-1)^n}{\sqrt n}$：\omterm{thm:b1:series:alternating}{交错级数}，且 $\frac{1}{\sqrt n}
\downarrow 0$，故收敛（见\cref{thm:b1:series:alternating}）；不
\omterm{thm:b1:series:absolute}{绝对收敛}（$\alpha = \frac12 \leq 1$）。

$\sum \frac{(-1)^n}{n + (-1)^n}$：数列 $\frac{1}{n +
(-1)^n}$ \emph{不是}递减的（$\frac{1}{n+1}$ 与
$\frac{1}{n}$ 交替出现），因此不能直接使用判别法。
展开：
\[
\frac{(-1)^n}{n + (-1)^n}
= \frac{(-1)^n}{n}\cdot\frac{1}{1 + \frac{(-1)^n}{n}}
= \frac{(-1)^n}{n} - \frac{1}{n^2} +
O\Bigl(\frac{1}{n^3}\Bigr):
\]
第一个\omterm{def:b1:series:def}{级数}收敛（\omterm{thm:b1:series:alternating}{交错级数}），$\sum \frac{1}{n^2}$
收敛，$O(n^{-3})$ 所对应的\omterm{def:b1:series:def}{级数}\omterm{thm:b1:series:absolute}{绝对收敛}：三个收敛\omterm{def:b1:series:def}{级数}之和
仍收敛。

$\sin\bigl(\pi\sqrt{n^2+1}\bigr)$：记 $\sqrt{n^2 + 1} = n +
\frac{1}{2n} + \varepsilon_n$，其中 $\varepsilon_n = O(n^{-3})$；
于是由 $\pi$ 周期性（对 $\sin$ 而言至多相差符号），
\[
\sin\bigl(\pi\sqrt{n^2+1}\bigr)
= (-1)^n \sin\Bigl(\frac{\pi}{2n} + \pi\varepsilon_n\Bigr) .
\]
记 $\theta_n = \frac{\pi}{2n} + \pi\varepsilon_n$，$a_n =
\sin\theta_n$。当 $n$ 充分大时，$\theta_n \in \intoo{0}{\frac\pi2}$
且
\[
\theta_n - \theta_{n+1} = \frac{\pi}{2n(n+1)} +
\pi(\varepsilon_n - \varepsilon_{n+1})
= \frac{\pi}{2n^2} + O\Bigl(\frac{1}{n^3}\Bigr) > 0
\]
从某项起 $(\theta_n)$ 递减并趋于 $0$；由于 $\sin$ 在
$\intcc{0}{\frac\pi2}$ 上递增，故 $(a_n)$ 也递减趋于
$0$。由\omterm{thm:b1:series:alternating}{交错级数}判别法可知\omterm{def:b1:series:def}{级数}收敛——但非\omterm{thm:b1:series:absolute}{绝对收敛}，
因为 $a_n \sim \frac{\pi}{2n}$。
\end{solution}

\begin{solution}{exo:b1:series:6}
$f(t) = \frac{1}{t(\ln t)^\alpha}$ 为正、\omterm{def:b1:continuity:continuous}{连续}、
在 $\intco{2}{+\infty}$ 上递减。作代换 $u = \ln t$：
\[
\int_2^x \frac{\dd t}{t(\ln t)^\alpha}
= \int_{\ln 2}^{\ln x} \frac{\dd u}{u^\alpha},
\]
当 $x \to \infty$ 时有界，当且仅当 $\alpha > 1$
（同\cref{thm:b1:series:riemann}中的计算）。由\omterm{thm:b1:integration:def}{积分}判别法
比较可得：\omterm{def:b1:series:def}{级数}收敛当且仅当 $\alpha > 1$。（这类\emph{贝特朗型}
\omterm{def:b1:series:def}{级数}显示了收敛\omterm{def:b1:topology:closure}{边界}的精细程度：$n\ln n$
发散，$n(\ln n)^{1.01}$ 收敛。）
\end{solution}

\begin{solution}{exo:b1:series:7}
由\cref{exo:b1:derivative:3}的切线界经展开改写：
取 $x = \frac1n \in \intoc{0}{1}$，对 $\ln(1+x)$ 取泰勒–拉格朗日公式
（$1$ 阶），得 $\ln(1 + x) = x - \frac{x^2}{2(1 +
c)^2}$，其中 $c \in \intoo{0}{x}$，于是
\[
0 \leq u_n = \frac1n - \ln\Bigl(1 + \frac1n\Bigr) \leq
\frac{1}{2n^2} .
\]
与\omterm{thm:b1:series:riemann}{黎曼级数}比较可知 $\sum u_n$ 收敛。其部分和
的对数项裂项相消：
\[
\sum_{n=1}^{N} u_n = H_N - \ln(N+1)
\]
（由于 $\sum_{n\leq N} \ln\frac{n+1}{n} = \ln(N+1)$）。故 $H_N -
\ln(N+1)$ 收敛；加上 $\ln\frac{N+1}{N} \to 0$ 后，数列
$H_N - \ln N$ 收敛，其极限为 $\gamma \approx 0.5772$。
\end{solution}

\begin{solution}{exo:b1:series:8}
$\dfrac{1}{n(n+2)} = \dfrac{1/2}{n} - \dfrac{1/2}{n+2}$，部分和
以间隔 $2$ 相消，
\[
S_N = \frac12\Bigl(1 + \frac12 - \frac{1}{N+1} -
\frac{1}{N+2}\Bigr) \longrightarrow \frac34 .
\]

$\sum \frac{n}{2^n}$：对 $\abs x < 1$，将有限几何和逐项求导，
再取极限（本题各\omterm{def:b1:series:def}{级数}均\omterm{thm:b1:series:absolute}{绝对收敛}，由\omterm{thm:b1:series:ratio}{比值判别法}）。
由 $\sum_{n\geq0} x^n = \frac{1}{1-x}$，
用部分和直接计算可得
\[
\sum_{n=1}^{N} n x^{n-1}
= \frac{1 - (N+1)x^N + N x^{N+1}}{(1 - x)^2}
\xrightarrow[N\to\infty]{} \frac{1}{(1-x)^2}
\quad (\abs x < 1),
\]
（\omterm{def:b1:topology:closure}{边界}项 $N x^N \to 0$）。取 $x = \frac12$：
$\sum_{n\geq1} n\bigl(\frac12\bigr)^{n-1} = 4$，故 $\sum_{n\geq0}
\frac{n}{2^n} = \frac12 \times 4 = 2$。
\end{solution}

\begin{solution}{exo:b1:series:9}
将 $\sum u_n$ 的项按 $2$ 的幂次分组：
\omterm{def:b1:reals:bounds}{上界}估计：当 $2^k \leq n < 2^{k+1}$ 时，共有 $2^k$ 项，
每项均 $\leq u_{2^k}$：
\[
\sum_{n=1}^{2^{K+1}-1} u_n
= \sum_{k=0}^{K} \sum_{n=2^k}^{2^{k+1}-1} u_n
\leq \sum_{k=0}^{K} 2^k u_{2^k} .
\]
下界估计：同一组中每项均 $\geq u_{2^{k+1}}$，
故 $\sum_{n=2^k}^{2^{k+1}-1} u_n \geq 2^k u_{2^{k+1}} = \frac12
\cdot 2^{k+1} u_{2^{k+1}}$，于是
\[
\sum_{n=1}^{2^{K+1}-1} u_n \geq \frac12 \sum_{k=1}^{K+1} 2^{k}
u_{2^{k}} .
\]
两个部分和的比较均可双向使用（各项非负，
由\cref{thm:b1:series:positive}）：两\omterm{def:b1:series:def}{级数}敛散性相同。

\omterm{thm:b1:series:riemann}{黎曼级数}：取 $u_n = n^{-\alpha}$，则 $2^k u_{2^k} =
2^{k(1-\alpha)}$，为\omterm{ex:b1:series:geometric}{等比级数}，当且仅当 $2^{1 - \alpha}
< 1$，即 $\alpha > 1$ 时收敛。贝特朗
（见\cref{exo:b1:series:6}）：取 $u_n = \frac{1}{n(\ln n)^\alpha}$，则
$2^k u_{2^k} = \frac{1}{(k\ln 2)^\alpha}$，为关于 $k$ 的\omterm{thm:b1:series:riemann}{黎曼级数}：
当且仅当 $\alpha > 1$ 时收敛。
\end{solution}

\begin{solution}{exo:b1:series:10}
公比为 $-t^2$ 的有限等比恒等式：
\[
\frac{1}{1 + t^2} = \sum_{k=0}^{n-1} (-1)^k t^{2k}
+ \frac{(-1)^n t^{2n}}{1 + t^2} .
\]
在 $\intcc{0}{1}$ 上\omterm{thm:b1:integration:def}{积分}（左边的\omterm{thm:b1:integration:def}{积分}为 $\arctan
1 = \frac\pi4$，见\cref{prop:b1:functions:arcderiv}）：
\[
\frac{\pi}{4} = \sum_{k=0}^{n-1} \frac{(-1)^k}{2k+1}
+ (-1)^n \int_0^1 \frac{t^{2n}}{1+t^2}\,\dd t,
\qquad
0 \leq \int_0^1 \frac{t^{2n}}{1+t^2}\,\dd t \leq \int_0^1 t^{2n}\dd
t = \frac{1}{2n+1} .
\]
令 $n \to \infty$ 即证得该公式，而\omterm{thm:b1:integration:def}{积分}的上述估计
恰为余项的界：对前
$N$ 项求和（即 $n = N + 1$ 项）后，$\abs{R_N} \leq \frac{1}{2N + 3}$。
\end{solution}

\begin{solution}{exo:b1:series:11}
$n^{1/n} = \eu^{\frac{\ln n}{n}} \to \eu^0 = 1$
（见\cref{prop:b1:functions:powerrules}）。于是
\[
\frac{1}{n^{1 + 1/n}} = \frac{1}{n}\,\eu^{-\frac{\ln n}{n}}
\sim \frac{1}{n} ,
\]
由等价判别法（见\cref{thm:b1:series:positive}）
与发散的调和\omterm{def:b1:series:def}{级数}比较：\emph{发散}，
虽然每个指数 $1 + \frac1n$ 都大于 $1$。黎曼判别法
针对\emph{固定}的指数 $\alpha$；
指数滑落至 $1$ 时可能丧失全部余量，本例正是如此。
\end{solution}

\begin{solution}{exo:b1:series:12}
设 $\varepsilon > 0$。由收敛\omterm{def:b1:series:def}{级数}的柯西准则
（对部分和应用\cref{thm:b1:seq:complete}），
存在 $N$，使得 $\sum_{k=n+1}^{2n} u_k \leq
\varepsilon$ 对 $n \geq N$ 成立。由单调性，这 $n$ 项中每项均
$\geq u_{2n}$：
\[
n\,u_{2n} \leq \sum_{k=n+1}^{2n} u_k \leq \varepsilon
\quad\Longrightarrow\quad 2n\,u_{2n} \leq 2\varepsilon ,
\]
对奇数指标，$(2n+1)\,u_{2n+1} \leq (2n+1)\,u_{2n} \leq
2\bigl(2n\,u_{2n}\bigr) \leq 4\varepsilon$，当 $n \geq N$ 时，
两种奇偶性下均有 $n u_n \to 0$。

逆命题不成立：$u_n = \frac{1}{n\ln n}$ 满足 $n u_n =
\frac{1}{\ln n} \to 0$，但\omterm{def:b1:series:def}{级数}发散
（见\cref{exo:b1:series:6}，$\alpha = 1$）。单调性不可去掉：
设 $u_n = \frac1n$，当 $n$ 为完全平方数时，$u_n =
2^{-n}$，其余情形：\omterm{def:b1:series:def}{级数}收敛（平方项之和形如
$\sum \frac{1}{k^2}$，其余项为几何\omterm{def:b1:series:def}{级数}），但在平方数处 $n u_n =
1$。
\end{solution}

\begin{solution}{pb:b1:series:1}
\textbf{1.} $a_{n+1} - a_n = \frac{1}{n+1} - \ln\frac{n+1}{n}
\leq 0$，因为 $\ln(1 + \frac1n) \geq \frac{1/n}{1 + 1/n} =
\frac{1}{n+1}$；而 $b_{n+1} - b_n = \frac{1}{n+1} -
\ln\frac{n+2}{n+1} \geq 0$，因为 $\ln(1 + \frac{1}{n+1}) \leq
\frac{1}{n+1}$。二者的间距 $a_n - b_n = \ln(1 + \frac1n) \to 0$：
相邻（见\cref{thm:b1:seq:adjacent}），共同极限 $\lim
a_n = \gamma$（见\cref{exo:b1:series:7}），且 $b_n \leq \gamma
\leq a_n$。

\textbf{2.} $b_{10} = 2.928968 - \ln 11 = 0.5311$，且 $a_{10} =
2.928968 - \ln 10 = 0.6264$，故 $\gamma \in
\intcc{0.5311}{0.6264}$。间距为 $\ln 1.1 \approx 0.095$，且
按 $\frac1n$ 的速度缩小：精确到小数点后四位（$\text{间距} \leq
10^{-4}$）需 $n \approx 10^4$——\omterm{prop:b1:reals:intervals}{区间}是
正确的，但很慢。

\textbf{3.} 裂项相消，$a_n - a_{m+1} = \sum_{k=n}^{m} w_k$，并
令 $m \to \infty$：$a_n - \gamma = \sum_{k\geq n} w_k$
（部分和的极限）。此外
\[
w_k = \int_k^{k+1} \frac{\dd t}{t} - \frac{1}{k+1}
= \int_k^{k+1} \Bigl(\frac1t - \frac{1}{k+1}\Bigr)\dd t
= \int_k^{k+1} \frac{k + 1 - t}{t\,(k+1)}\,\dd t .
\]
在 $\intcc{k}{k+1}$ 上：$\frac{1}{(k+1)^2} \leq \frac{1}{t(k+1)}
\leq \frac{1}{k(k+1)}$，且 $\int_k^{k+1}(k + 1 - t)\dd t =
\frac12$，因此 $\frac{1}{2(k+1)^2} \leq w_k \leq
\frac{1}{2k(k+1)}$。

\textbf{4.} \omterm{def:b1:reals:bounds}{上界}：$\sum_{k \geq n} \frac{1}{2k(k+1)} =
\frac12\sum_{k\geq n}\bigl(\frac1k - \frac{1}{k+1}\bigr) =
\frac{1}{2n}$（裂项相消）。下界：$\frac{1}{2(k+1)^2} \geq
\frac{1}{2(k+1)(k+2)}$，其和裂项相消得
$\frac{1}{2(n+1)}$。结合第 3 问：
\[
\frac{1}{2(n+1)} \leq H_n - \ln n - \gamma \leq \frac{1}{2n} .
\]

\textbf{5.} 减去 $\frac{1}{2n}$：$\gamma_n - \gamma \in
\intcc{\frac{1}{2(n+1)} - \frac{1}{2n}}{0} =
\intcc{-\frac{1}{2n(n+1)}}{0}$：修正后的估计
精确到 $O\bigl(\frac{1}{n^2}\bigr)$ 阶，且始终从下方逼近。

\textbf{6.} $\gamma_{100} = 5.1873775 - \ln 100 - 0.005 =
0.5772073$，并有 $0 \leq \gamma - \gamma_{100} \leq
\frac{1}{20200} < 5\cdot10^{-5}$，因此 $0.577207 \leq \gamma
\leq 0.577257$，即 $\gamma = 0.5772 \pm 5\cdot10^{-5}$（真值
为 $0.5772156\dots$）——由一百项即得四位可靠小数，
相比第 2 问所需的一万项。

\textbf{7.} 第一部分：
\[
H_{2m} - H_m = \Bigl(\ln 2m + \gamma + \frac{1}{4m}\Bigr)
- \Bigl(\ln m + \gamma + \frac{1}{2m}\Bigr) +
O\Bigl(\frac{1}{m^2}\Bigr)
= \ln 2 - \frac{1}{4m} + O\Bigl(\frac{1}{m^2}\Bigr).
\]
第二部分：不超过 $2m$ 的偶数倒数之和为 $\frac12 H_m$，故
$\sum_{j=1}^{m}\frac{1}{2j-1} = H_{2m} - \frac12 H_m = \frac12
\ln m + \ln 2 + \frac\gamma2 + o(1)$，且 $\sum_{j=1}^m
\frac{1}{2j} = \frac12\ln m + \frac\gamma2 + o(1)$。

\textbf{8.} \omterm{def:b1:continuity:continuous}{连续}两次分出偶数项：
$\sum_{k=1}^{2m}\frac{(-1)^{k-1}}{k} = H_{2m} - 2\cdot\frac12
H_m = H_{2m} - H_m$。由第 7 问，此式等于 $\ln 2 -
\frac{1}{4m} + O(m^{-2})$：故和为 $\ln 2$
（再次见\cref{ex:b1:series:ln2}）\emph{连同}其收敛速度。

\textbf{9.} 当 $N = 2m$ 时：$S - S_N = \frac{1}{4m} + O(m^{-2}) =
\frac{1}{2N} + O(N^{-2})$。当 $N = 2m + 1$ 时：$S_{N} = S_{2m} +
\frac{1}{2m+1}$，于是
\[
S - S_N = \Bigl(\frac{1}{4m} - \frac{1}{2m+1}\Bigr) +
O\Bigl(\frac{1}{m^2}\Bigr)
= -\frac{1}{4m} + O\Bigl(\frac{1}{m^2}\Bigr)
= -\frac{1}{2N} + O\Bigl(\frac{1}{N^2}\Bigr).
\]
两种情形下均有 $S - S_N \sim \frac{(-1)^N}{2N}$：恰为
最坏界 $a_{N+1}$ 的一半，且符号已知并交替。

\textbf{10.} 取平均可消去振荡的主项：
\[
S - \tilde S_N = \frac{(S - S_N) + (S - S_{N+1})}{2}
= \frac{(-1)^N}{2}\Bigl(\frac{1}{2N} - \frac{1}{2(N+1)}\Bigr)
+ O\Bigl(\frac{1}{N^2}\Bigr) = O\Bigl(\frac{1}{N^2}\Bigr).
\]
数值上：$S_{10} = 0.645635$，$S_{11} = 0.736544$，$\tilde
S_{10} = 0.691089$，与 $\ln 2 = 0.693147$：误差由
$4.8\cdot10^{-2}$ 降至 $2.1\cdot10^{-3}$——仅一次加法，
效果提高二十倍。

\textbf{11.} 交错误差每一步都变号，
因此相邻部分和夹住极限，其平均
消去一阶项；\omterm{prop:b1:reals:intervals}{区间}误差 $H_n - \ln n -
\gamma \approx \frac{1}{2n}$ 符号恒定，故沿
$n$ 的任何平均都无法消去它。对两个\emph{\omterm{prop:b1:reals:intervals}{区间}}取平均确有帮助：
$\frac{a_n + b_n}{2} = H_n - \ln\sqrt{n(n+1)}$，且由于
$\ln\sqrt{n(n+1)} = \ln\bigl(n + \frac12\bigr) + O(n^{-2})$,
\[
H_n - \ln\Bigl(n + \frac12\Bigr)
= \Bigl(H_n - \ln n - \frac{1}{2n}\Bigr) + \frac{1}{8n^2} +
O\Bigl(\frac{1}{n^3}\Bigr) = \gamma + O\Bigl(\frac{1}{n^2}\Bigr)
\]
（第 5 问及 $\ln(1 + \frac{1}{2n}) = \frac{1}{2n} -
\frac{1}{8n^2} + O(n^{-3})$）。取 $n = 10$ 检验：$H_{10} -
\ln 10.5 = 0.57759$，已在
$4\cdot10^{-4}$ 之内逼近 $\gamma$。

\textbf{12.} $\sum_{j\leq m} \frac{1}{2j-1} \geq \sum_{j \leq m}
\frac{1}{2j} = \frac12 H_m \to +\infty$：交错调和\omterm{def:b1:series:def}{级数}的正部与
负部都发散。

\textbf{13.} 记 $u_n^\pm = \frac{\abs{u_n} \pm u_n}{2} \geq
0$，于是 $u_n = u_n^+ - u_n^-$ 且 $\abs{u_n} = u_n^+ + u_n^-$。若
$\sum u_n^+$ 收敛，则 $\sum u_n^- = \sum (u_n^+ - u_n)$
亦收敛（收敛\omterm{def:b1:series:def}{级数}之差），从而 $\sum
\abs{u_n}$ 亦收敛，与条件收敛矛盾。由对称性，
$\sum u_n^\pm$ 均发散（趋于 $+\infty$）：
即正、负质量各为无穷的储库。

\textbf{14.} 设 $S = \sum u_n$，$\varepsilon > 0$，且 $N$ 满足
$\sum_{n > N} \abs{u_n} \leq \varepsilon$（$\sum\abs{u_n}$ 的柯西
准则）。设 $M_0$ 足够大，使 $\sigma(\{0,
\dots, M_0\}) \supseteq \{0, \dots, N\}$。当 $M \geq M_0$ 时，
$\sum_{m \leq M} u_{\sigma(m)} - \sum_{n \leq N} u_n$ 之差
是若干\emph{互异}项 $u_n$（$n > N$）的有限和，
从而其绝对值 $\leq \varepsilon$；又 $\abs{S -
\sum_{n\leq N} u_n} \leq \varepsilon$ 亦成立。于是重排后的
部分和以 $2\varepsilon$ 为误差最终逼近 $S$：
$\sum u_{\sigma(n)} = S$。\omterm{thm:b1:series:absolute}{绝对收敛}
具有抗重排性。

\textbf{15.} 贪心过程的每一阶段都在
有限多项之后结束，因为仅剩的正项
（相应地，负项）单独的部分和发散
（第 12 问）：累加和必定最终越过 $t$。该
过程因此交替进行无穷多个有限阶段，
依次取用正项与负项，且各自保持原有顺序：
每一项恰好用一次——这便是一个重排。
首次越过之后，相邻两次越过之间的
部分和单调地趋向 $t$，而在越过处
超出量至多为刚加入的那一项；由于
第 $j$ 次越过时所用各项在其类中的指标至少为 $j$，
故这些超出量趋于 $0$。于是部分和
收敛于 $t$：每个实数都是某个
重排之和。

\textbf{16.} 正位依次填入 $\frac{1}{2j-1}$，其中
$j = 1, 2, \dots$ 依次取值，负位依次填入 $\frac{1}{2j}$，
故交错调和\omterm{def:b1:series:def}{级数}的每一项都出现，
且恰好出现一次。取 $(p, q) = (1, 2)$，各块为
$\bigl(1, -\frac12, -\frac14\bigr)$, $\bigl(\frac13, -\frac16,
-\frac18\bigr)$, $\bigl(\frac15, -\frac1{10},
-\frac1{12}\bigr)$，\dots\ ——即所示\omterm{def:b1:series:def}{级数}。

\textbf{17.} 由于 $\frac{1}{4k-2} = \frac12\cdot\frac{1}{2k-1}$：
\[
\frac{1}{2k-1} - \frac{1}{4k-2} - \frac{1}{4k}
= \frac12\,\frac{1}{2k-1} - \frac12\,\frac{1}{2k}
= \frac12\Bigl(\frac{1}{2k-1} - \frac{1}{2k}\Bigr).
\]
对 $k = 1, \dots, K$ 求和：$T_{3K} = \frac12
\sum_{k=1}^{K}\bigl(\frac{1}{2k-1} - \frac{1}{2k}\bigr) =
\frac12 S_{2K}$：每逢第三个部分和，
重排\omterm{def:b1:series:def}{级数}\emph{恰好}为原\omterm{def:b1:series:def}{级数}的一半。

\textbf{18.} 完成 $K$ 个完整块之后，重排后的部分和
为 $\sum_{j=1}^{pK}\frac{1}{2j-1} -
\sum_{j=1}^{qK}\frac{1}{2j}$，由第 7 问算得：
\[
\Bigl(\frac{\ln(pK)}{2} + \ln 2 + \frac\gamma2\Bigr)
- \Bigl(\frac{\ln(qK)}{2} + \frac\gamma2\Bigr) + o(1)
= \ln 2 + \frac12\ln\frac pq + o(1) :
\]
$\gamma$ 相消，$\ln K$ 相消，比值
$\frac pq$ 保留下来。

\textbf{19.} 第 $K + 1$ 块内的部分和，与
第 $K$ 块之和相差至多 $p + q$ 项，每项绝对值
$\leq \frac{1}{2qK}$ 左右，因此相差 $O\bigl(\frac1K\bigr) \to 0$：
部分和的整个序列具有相同极限 $\ln 2 +
\frac12\ln\frac pq$。对 $(1, 2)$，$\ln 2 + \frac12\ln\frac12 =
\frac{\ln 2}{2} = 0.34657\dots$，且 $T_9 = 0.30833$
正逼近它：由第 17 问，$T_{3K} = \frac12 S_{2K}$
收敛，且其误差恰为交错调和\omterm{def:b1:series:def}{级数}误差的一半。
项相同，和减半。

\textbf{20.} $(1,1)$：$\ln 2 + \frac12\ln 1 = \ln 2$ ——
原有顺序，自洽。$(2,1)$：$\frac32\ln 2 \approx
1.0397$。$(p,q)$ 的选取恰好给出可数\omterm{def:b1:topology:dense}{稠密}的
族 $\ln 2 + \frac12\ln r$，$r \in \Q_{>0}$；黎曼的贪心
方案（第 15 问）可达\emph{每一个}实数。结构换来
公式，贪心换来全面性。

\textbf{21.} 部分和裂项相消：$\sum_{k=1}^{N}
\bigl(\frac1k - \ln\frac{k+1}{k}\bigr) = H_N - \ln(N+1) = b_N
\to \gamma$：\cref{exo:b1:series:7} 中的\omterm{def:b1:series:def}{级数}恰好
和为\omterm{pb:b1:series:1}{欧拉常数}。

\textbf{22.} 对 $t = 1$ 作贪心选取：第一个正项使
和恰好为 $1$，并未越过，故须再取第二个正项
方可越过：$1, \frac13$（和 $1.3333 > 1$），然后
$-\frac12$（$0.8333$），$\frac15$（$1.0333$），$-\frac14$
（$0.7833$），$\frac17, \frac19$（$1.0373$），$-\frac16$
（$0.8706$），$\frac1{11}, \frac1{13}$（$1.0384$），$-\frac18$
（$0.9134$），$\frac1{15}$（$0.9801$），\dots\ ——部分和围绕
$1$ 附近振荡，振幅愈来愈小，
每个周期现需两个正项，因为负项的绝对值更大。

\textbf{23.} 构造分块：在第 $j$ 阶段，接续添加足够多的未用
正项，使部分和至少增加 $1$
（可以做到：剩余正项的和发散，
第 12 问），然后接续添加单个负项 $-\frac{1}
{2j}$。每个正项终将被用到（每一阶段至少用一个），
每个负项也同样（每阶段一个）：这是一个
重排。每一阶段使和的改变量 $\geq 1 -
\frac{1}{2j} \geq \frac12$，故部分和
超过 $\frac{j}{2}$（第 $j$ 阶段之后），而阶段内的增量
除最后一个外均为正，最后一个增量以 $\frac{1}{2j} \to 0$ 为界：
于是发散到 $+\infty$。

\textbf{24.} 由该规律，$H_N \geq 20 \iff \ln N \geq 20 -
\gamma - \frac{1}{2N} + O(N^{-2})$：阈值 $N^*$
满足 $\ln N^* = 20 - \gamma + o(1)$，即 $N^* =
\eu^{20 - \gamma}(1 + o(1)) \approx \eu^{19.4228} \approx
2.72\cdot10^{8}$ ——落在
$\intcc{1.8\cdot10^8}{4.9\cdot10^8}$ 所界定的粗略\omterm{prop:b1:reals:intervals}{区间}之内，
见 \cref{ex:b1:series:harmonicstack}，并由 $\gamma$ 所确定。

\textbf{25.} (i) 在 $H_n = \ln n + \gamma + \frac{1}{2n} +
O(n^{-2})$ 中，$\ln n$ 即\omterm{thm:b1:integration:def}{积分}，$\gamma$ 是用\omterm{thm:b1:integration:def}{积分}代替和式所付出的代价
（一个真正新的分析学常数），
而 $\frac{1}{2n}$ 为第一修正项——梯形法则的
影子。(ii) 条件收敛依赖于
两个无穷储库之间的相消（第 13 问），因此
重排会改变两储库的权重；\omterm{thm:b1:series:absolute}{绝对收敛}的
总质量有限，而第 14 问的尾部估计
与次序无关。(iii) $(p,q)$ 的和经过\emph{计算}得出：
$\gamma$ 规律把每个重排部分和化为 $\ln 2 +
\frac12\ln\frac pq + o(1)$，$\gamma$ 本身相消——
一个渐近计算练习，而非抽象论证。
(iv) 下一步：幂\omterm{def:b1:series:def}{级数}的乘积与无条件可和性，
见本科二年级分册；以及本科三年级分册的周末问题
关于斯特林公式，其中和与\omterm{thm:b1:integration:def}{积分}的
核算再进一阶，便得出 $\sqrt{2\pi}$
本身。
\end{solution}
